\section{The ready-cohort model}
\label{sec:model}

\subsection{Events and the hardware threshold}

Let $E$ be a set of control events. Event $i$ has release time $t_i$, latest
admissible launch time $d_i\geq t_i$, and route $r_i$. A route in the
mathematical model means that one implementation can process the grouped events
without changing their declared transition semantics. A trace label does not
establish that condition by itself.

For route $r$ and hardware/runtime configuration $h$, define
$K_r(h,v,H)$ as the start of a measured safe suffix: every tested cohort size
from $K_r$ through a declared maximum beats a named admissible baseline. The
model extrapolates that monotone threshold and has no upper batch limit. The
arguments record state visibility $v$ and observation-free horizon $H$ because
transfer, synchronization, and temporal fusion can move the crossover. We
write $K_r$ when the configuration is fixed.

A batch $(\tau,B)$ is feasible when

\begin{equation}
  |B|\geq K_r,\quad r_i=r\ \forall i\in B,\quad
  t_i\leq\tau\leq d_i\ \forall i\in B.
  \label{eq:feasible-batch}
\end{equation}

A schedule contains feasible batches and assigns each event at most once. The
model allows batches to have more than $K_r$ events. It has no maximum batch
size, service time, or device-capacity constraint.

\subsection{Four shares with different meanings}

Let $\pi$ be a frozen partition of time into non-overlapping windows. Each
event belongs to a bucket defined by its window and route. If bucket $b$ has
$n_b$ events, fixed-partition eligibility is

\begin{equation}
  F(\pi,K)=\frac{1}{|E|}\sum_b n_b\,\mathbf{1}[n_b\geq K_{r(b)}].
  \label{eq:fixed}
\end{equation}

$F$ is exact only for schedulers constrained to $\pi$. Sliding deadlines can
join events across a fixed boundary, so $F$ is not a universal ceiling.

Let $\pstar$ be the maximum number of assigned events in any feasible offline
schedule, divided by $|E|$. It uses future knowledge and therefore measures
opportunity rather than online performance.

For route $r$, define the active count at time $\tau$ as

\begin{equation}
  Q_r(\tau)=|\{j:r_j=r,\ t_j\leq\tau\leq d_j\}|.
\end{equation}

Event $i$ is locally eligible if $Q_{r_i}(\tau)\geq K_{r_i}$ at some
$\tau\in[t_i,d_i]$. The local-overlap share is

\begin{equation}
  U=\frac{1}{|E|}\sum_i \mathbf{1}\!\left[
  \max_{\tau\in[t_i,d_i]}Q_{r_i}(\tau)\geq K_{r_i}\right].
  \label{eq:upper}
\end{equation}

$U$ may count incompatible opportunities that reuse the same events. It is an
upper bound, not necessarily an achievable schedule.

Finally, $A$ is the share accelerated within deadline by a specified finite
online runtime, including its queue, service, fallback, and capacity rules. To
compare $A$ with $\pstar$, each accelerated batch must obey the same route
grouping, threshold $K_r$, and launch deadline as the offline model; finite
service and capacity may impose stricter constraints. Its denominator is the
identical event set $E$ and retained horizon used by $\pstar$; late, missing,
failed, and fallback events remain in that denominator. The current paper does
not measure $A$.

\begin{proposition}[Boundary ordering]
\label{prop:ordering}
If every batch admitted by the frozen partition is admissible under the event
deadlines, and the batches counted by $A$ obey the same route, threshold,
deadline, and event-set constraints as the offline model, then
\begin{equation}
  F\leq\pstar\leq U,\qquad A\leq\pstar.
  \label{eq:ordering}
\end{equation}
\end{proposition}

\begin{proof}
Every bucket with $n_b\geq K_{r(b)}$ is a feasible batch, and buckets are
disjoint, so the fixed schedule attains $F$. Every event in any feasible batch
overlaps at least $K_{r_i}$ same-route intervals at that batch's launch time,
so it is counted by $U$. An online schedule is one member of the offline
feasible set.
\end{proof}

\subsection{Compatibility has a measurable tax}

Suppose grouping $g_f$ refines $g_c$, so every fine bucket lies inside one
coarse bucket. For fixed $\pi$ and $K$,

\begin{equation}
  F(\pi,K,g_f)\leq F(\pi,K,g_c).
  \label{eq:grouping}
\end{equation}

This inequality measures fragmentation under a declared grouping. It does not
make coarse fusion admissible.
A pooled kernel may change control flow, memory access, actions, or numerical
trajectory. The route-key proxy is therefore the narrowest grouping in the
trace study.

\begin{table}[t]
  \centering
  \footnotesize
  \begin{tabularx}{\columnwidth}{@{}llX@{}}
    \toprule
    Symbol & Source & Meaning \\
    \midrule
    $K_r$ & hardware & route-specific profitable cohort \\
    $F$ & workload, policy & fixed-partition eligible share \\
    $\pstar$ & workload, offline & exact schedulable share \\
    $U$ & workload, local & necessary-overlap upper bound \\
    $A$ & online runtime & achieved accelerated share \\
    \bottomrule
  \end{tabularx}
  \caption{The quantities cannot be substituted for one another. In
  particular, $K=256$ in the trace sweep is not a measured $K_r$ for the
  resident-policy mechanism.}
  \label{tab:notation}
\end{table}
